\section{Conclusions}

\begin{frame}{What is established}
  \bookref{Synthesis of Chapters 1--6}
  \begin{enumerate}
    \item \textbf{Finite invariant:}
      $I(S)=\log(n!/|\Aut(S)|)$ is exactly the logarithm of a relabeling orbit.
    \item \textbf{Structural calculus:}
      bounds, rigidity, refinement monotonicity, and preservation results follow
      from subgroup structure.
    \item \textbf{Statistical role:}
      $e^{I(S)}$ is the isomorphism-class multiplicity in relational partition
      functions.
    \item \textbf{Explicit transition:}
      the refined two-block model has an exact antisymmetric instability at
      $\beta_c=2/J$ and a positive quartic Landau coefficient.
    \item \textbf{Extension principle:}
      operations are functional relations, so the same finite theory applies to
      algebraic structures.
    \item \textbf{Reference principle:}
      physical deficits compare $H$ with $G\leq\Aut(B)$, not with bare-set
      permutations by default.
  \end{enumerate}
\end{frame}

\begin{frame}{The next mathematical problems}
  \bookref{Chapter 3, Section 12; Chapters 5--6; Appendix C}
  \begin{itemize}
    \item Prove uniform remainder estimates and a genuine local bifurcation
      theorem for the exact refined free energy.
    \item Relate microscopic automorphism loss to graphon or kernel variational
      limits without conflating orbit entropy and block entropy.
    \item Determine when symmetry-deficit observables have stable finite-size
      scaling and how their exponents transform under nonlinear proxies.
    \item Construct a relational renormalization map and identify its fixed
      points and relevant directions.
    \item Prove exhaustion-independent limits for finite-window deficits and
      compare their scaling with conventional susceptibilities.
    \item Extend formal verification from algebraic derivatives to analytic
      estimates and limiting arguments.
  \end{itemize}
\end{frame}

\begin{frame}{Take-home message}
  \bookref{Introduction and final synthesis}
  \begin{center}
    {\Large
      \textcolor{RelBlue}{Symmetry deficit turns the slogan}

      \vspace{0.5em}
      \textbf{``structure arises through symmetry breaking''}

      \vspace{0.5em}
      \textcolor{RelBlue}{into an exact finite invariant.}}
  \end{center}

  \vspace{1em}
  \begin{block}{Broader program}
    Orbit counting gives the finite foundation; variational models show how the
    invariant can diagnose organized phases. Physical applications specify a
    structured reference group already at finite size, then use index, rate,
    geometric, or measure diagnostics as the state space and limits require.
  \end{block}
\end{frame}

\begin{frame}[allowframebreaks]{Selected references}
  \scriptsize
  \bibliographystyle{plain}
  \bibliography{../references}
\end{frame}
